% !TeX program = xelatex
% Template author: Dr. Awanish Pratap Singh
% XeLaTeX/LuaLaTeX: Arial-capable release build.
% pdfLaTeX: Helvetica drafting fallback; not permitted in submission mode.
% Generic English DFG Research Grant project-description template.
% Baseline: DFG forms 54.01 [06/26] and 53.01 elan [09/25].
% Source set rechecked: 2026-07-31.
% Independent convenience template; the official DFG forms remain authoritative.
%
% Drafting:   \documentclass[current,guidance,english]{dfgproposal}
% Submission: \documentclass[current,submission,english]{dfgproposal}
% Submission mode stops if any \DFGInstruction or \DFGDecisionRequired remains.
\documentclass[current,guidance,english]{dfgproposal}
% Current-source matrix shared by the English and German templates.
% Maintainer: Dr. Awanish Pratap Singh.
% Verified against the official DFG forms page on 2026-07-31.
% Recheck and update this file before starting a future submission.

\DFGSetFormLabel{DFG form 53.01}
\DFGSetProjectDescriptionRevision{09/25}
\DFGSetProposalInstructionsRevision{06/26} % form 54.01
\DFGSetMainPageMaximum{17}
\DFGSetSupplementPageMaximum{8}

\newcommand{\DFGResearchGrantsProgrammeRevision}{03/25} % form 50.01
\newcommand{\DFGBasicModuleRevision}{09/24}
\newcommand{\DFGPublicationListRevision}{03/25}
\newcommand{\DFGCurriculumVitaeRevision}{07/25}
\newcommand{\DFGAnimalCostsRevision}{12/25}
\newcommand{\DFGCoreFacilitiesRevision}{01/26}
\newcommand{\DFGSequencingServicesRevision}{03/24}
\newcommand{\DFGEquityDiversityRevision}{03/25}


\DFGSetDocumentTitle{Project Description -- Project Proposals}
\DFGSetApplicantName{[First name, last name, city of each applicant]}
\DFGSetProjectTitle{[Project title]}
\DFGSetKeywords{}

\begin{document}
\DFGStartMainMatter
\DFGMakeProposalHeader
\DFGMainPageLimitNotice

% Sections 1--3 together: maximum 17 pages.
\DFGSection{1}{Starting Point}
\DFGStackedUnnumberedHeading{State of the art and preliminary work}

\DFGInstruction{For a new proposal, describe the current state of the art in direct relation to the project, identify the precise gap, explain the project's original contribution, and present the applicants' relevant preliminary work. Make the argument understandable without requiring reviewers to consult cited literature. Identify applicant work clearly and cite only publicly available work, except accepted manuscripts accompanied by the required evidence. For a renewal, replace this with the DFG-required report on previous work.}

\DFGSection{2}{Objectives and work programme}

\DFGStackedSubsection{2.1}{Anticipated total duration of the project}

\DFGInstruction{State the requested funding period and the anticipated total duration. For a renewal or a project extending across funding periods, distinguish the present request from the full planned duration.}

\DFGSubsection{2.2}{Objectives}

\DFGInstruction{State the central research question or hypothesis, the specific objectives, and the expected scientific advance. Explicitly indicate whether the expected results may have relevance beyond science, for example for science policy, technology, the economy, or society. Keep broader-impact claims proportionate to the proposed evidence.}

\DFGSubsection{2.3}{Work programme incl. proposed research methods}

\DFGUnnumberedHeading{Applicant: [last name, first name]}

\DFGInstruction{Describe the work programme separately for each applicant. Define work packages, methods, experimental or analytical units, primary outcomes, sample-size or power logic, statistical analyses, quality controls, milestones, dependencies, responsibilities, risks, alternatives, and decision rules. Include a schedule detailing all planned experiments and other major tasks. Explain what is available locally and why any external contribution is needed. Repeat the applicant heading where several applicants share responsibility.}

\DFGSubsection{2.4}{Handling of research data}

\DFGInstruction{Describe data types and volumes; collection or generation; documentation and metadata; quality assurance; storage, backup, access and security; legal or ethical constraints; responsibilities; retention; publication and reuse; repositories; software or environment preservation; and the participating institution's concrete contribution to data and information management. State project costs needed for these activities where applicable.}

\DFGSubsection{2.5}{Relevance of sex, gender and/or diversity in the research project}

\DFGInstruction{Explain whether and how sex, gender and/or other diversity dimensions are relevant to the research questions, methods, work programme, objectives, analysis, or implementation. Consider people, people affected by results, animals, human or animal samples, researchers, and other relevant dimensions. If a dimension is not relevant, give a concise scientific justification.}

\DFGSection{3}{Project- and subject-related list of publications}
\DFGStartReferences

\DFGInstruction{List only works cited in sections 1 and 2. Give complete bibliographic details, including full titles and a DOI or stable URL where available; do not include journal impact factors, h-indices or comparable metrics. There is no total numerical cap, but no more than ten of the applicants' own works may be highlighted across all applicants. Non-public work is not citable, except an accepted manuscript submitted with the editor's confirmation. Keep this section at least Arial 9 pt and follow DFG form 1.91.}

\DFGEndReferences

% Section 4 onward together: maximum 8 pages. This command starts page 1 of 8.
\DFGStartSupplement
\DFGSection{4}{Supplementary information on the research context}
\DFGSupplementPageLimitNotice

\DFGStackedSubsection{4.1}{Ethical and/or legal aspects of the project}

\DFGStackedSubsubsection{4.1.1}{General ethical aspects}

\DFGInstruction{Identify discipline-specific ethical or legal issues, foreseeable risks or harms, governance, approvals, responsibilities, and timing. State whether an ethics-committee statement is required. If no issue applies, say so accurately and briefly.}

\DFGSubsubsection{4.1.2}{Descriptions of proposed investigations on humans, human materials or identifiable data}

\DFGInstruction{Where applicable, describe participant or material selection, sample-size justification, potential risks and precautions, information and consent, data protection, and the status of ethics approval. Otherwise use \string\DFGNotApplicable{} with a project-specific reason.}

\DFGSubsubsection{4.1.3}{Descriptions of proposed investigations involving experiments on animals}

\DFGInstruction{Where applicable, justify the animal model and procedures; report species, sex, age or developmental stage, group structure, animal numbers and statistical design; explain reduction, refinement and replacement; describe welfare, severity, humane endpoints and authorisation status; and reconcile numbers with the work programme and budget. Otherwise use \string\DFGNotApplicable{}.}

\DFGSubsubsection{4.1.4}{Descriptions of projects involving genetic resources (or associated traditional knowledge) from a foreign country}

\DFGInstruction{Address access-and-benefit-sharing obligations and compliance where genetic resources or associated traditional knowledge from another country are involved. Otherwise use \string\DFGNotApplicable{}.}

\DFGSubsubsection{4.1.5}{Explanations regarding any possible safety-related aspects}

\DFGInstruction{Describe safety-relevant aspects, risk assessment, governance, responsible bodies, approvals, and mitigations, or state that the subsection is not applicable.}

\DFGDeepSubsubsection{4.1.5.1}{``Dual Use Research of Concern''; foreign trade law}

\DFGInstruction{Assess possible dual-use concerns, export-control or foreign-trade restrictions, sensitive knowledge transfer, and required consultations or licences. State not applicable only after a project-specific assessment.}

\DFGDeepSubsubsection{4.1.5.2}{Risks in international cooperation}

\DFGInstruction{Where international cooperation is planned, assess the risk-benefit balance for the research subject, partners and conditions, and state proportionate risk-reduction measures. Otherwise use \string\DFGNotApplicable{}.}

\DFGSubsubsection{4.1.6}{Considerations on aspects of ecological sustainability in the planning and implementation of the project}

\DFGInstruction{Provide a brief project-specific reflection on sustainability in the planned research processes. Preserve research quality and explain proportionate resource- or emission-reducing measures; identify any additional funds requested for a more sustainable approach.}

\DFGSubsection{4.2}{Employment status information}

\DFGInstruction{For each applicant, state last name, first name and employment status. For a fixed-term contract, include its duration and funding body.}

\DFGSubsection{4.3}{First-time proposal data}

\DFGInstruction{State the last name and first name of each applicant submitting a first proposal to the DFG, or use \string\DFGNotApplicable{}. Apply the programme-specific definition from current form 54.01.}

\DFGSubsection{4.4}{Composition of the project group}

\DFGInstruction{List only people who will work on the project but will not be paid from the requested project funds. Give name, academic title, employment status and funding type. Do not add written descriptions of individuals' diversity characteristics.}

\DFGSubsection{4.5}{Researchers in Germany with whom you have agreed to cooperate on this project}

\DFGInstruction{List co-investigators and other researchers in Germany who will collaborate without sharing applicant responsibility, and identify required cooperation statements. For clinical trials, include the responsible biometrician or statistician. Otherwise use \string\DFGNotApplicable{}.}

\DFGSubsection{4.6}{Researchers abroad with whom you have agreed to cooperate on this project}

\DFGInstruction{List foreign cooperation partners and classify the cooperation under the applicable DFG international measure or general international cooperation route. Add required commitments and elan classifications. Otherwise use \string\DFGNotApplicable{}.}

\DFGSubsection{4.7}{Researchers with whom you have collaborated scientifically within the past three years}

\DFGInstruction{Provide the complete conflict-of-interest information requested by the current DFG definition and reporting period. Do not replace this with a selected collaborator list.}

\DFGSubsection{4.8}{Project-relevant cooperation with commercial enterprises}

\DFGInstruction{Describe any project-relevant cooperation with a commercial enterprise, state-aid implications and required agreements or declarations. Otherwise use \string\DFGNotApplicable{}.}

\DFGSubsection{4.9}{Project-relevant participation in commercial enterprises}

\DFGInstruction{Disclose ownership, directorship or other participation in a commercial enterprise and explain any link to the project's subject or the enterprise's activities. Otherwise use \string\DFGNotApplicable{}.}

\DFGSubsection{4.10}{Scientific equipment}

\DFGInstruction{List larger instruments and major computing facilities available for the project, including location, access and relevant capacity. If requested instrumentation exists at the institution but is unavailable to the project, explain why.}

\DFGSubsection{4.11}{Other submissions}

\DFGInstruction{List previous or parallel submissions for this project and/or major instrumentation, with status and overlap, or use \string\DFGNotApplicable{}.}

\DFGSubsection{4.12}{Other information}

\DFGInstruction{Add only processing-relevant information not reported elsewhere. Check the current form 54.01 for any mandatory proposal-preparation disclosure and include an applicant-approved factual statement where required.}

\DFGSection{5}{Requested modules}

\DFGInstruction{Explain each requested item separately for each applicant, stating last name and first name. Reconcile every amount with elan, the work programme, schedules, cohort or sample ledgers, and quotations. Keep the official module headings; use \string\DFGNotRequested{} for modules not requested.}

\DFGStackedSubsection{5.1}{Basic Module}

\DFGStackedSubsubsection{5.1.1}{Funding for Staff}

\DFGInstruction{For each requested position, state applicant, category, duration, percentage, project tasks, timing and scientific justification. Use current DFG personnel rates in elan.}

\DFGSubsubsection{5.1.2}{Direct Project Costs}

\DFGStackedDeepSubsubsection{5.1.2.1}{Equipment up to €10,000, Software and Consumables}

\DFGInstruction{Itemise quantities, unit costs, timing, applicant and work-package purpose. Confirm non-availability from basic institutional resources and remove double counting.}

\DFGDeepSubsubsection{5.1.2.2}{Travel}

\DFGInstruction{State purpose, destination or calculation basis, participants, timing and scientific justification, or use \string\DFGNotRequested{}.}

\DFGDeepSubsubsection{5.1.2.3}{Visiting Researchers (excluding Mercator Fellows)}

\DFGInstruction{State visitor, purpose, duration, calculation and scientific contribution, or use \string\DFGNotRequested{}.}

\DFGDeepSubsubsection{5.1.2.4}{Expenses for Laboratory Animals}

\DFGInstruction{Reconcile species, acquisition or breeding, numbers, housing weeks, unit rates and totals with sections 2.3 and 4.1.3. Apply current form 55.03, request only project-specific additional expenditure, and exclude institutional baseline and authority fees.}

\DFGDeepSubsubsection{5.1.2.5}{Other Costs}

\DFGInstruction{Itemise other project costs, including eligible contracts or core-facility usage. For internal core facilities or instrumentation-related usage costs, apply current form 55.04 and attach an itemised facility estimate when the applicable threshold is exceeded.}

\DFGDeepSubsubsection{5.1.2.6}{Project-related Publication Expenses}

\DFGInstruction{State the amount and calculation, avoid duplication with institutional coverage, and verify the current programme ceiling at submission, or use \string\DFGNotRequested{}.}

\DFGSubsubsection{5.1.3}{Instrumentation}

\DFGStackedDeepSubsubsection{5.1.3.1}{Equipment Exceeding €10,000}

\DFGInstruction{Identify the item, gross price, applicant, utilisation, scientific need, available alternatives, market comparison, host follow-up costs and required quotations, or use \string\DFGNotRequested{}.}

\DFGDeepSubsubsection{5.1.3.2}{Major Instrumentation Exceeding €50,000}

\DFGInstruction{Identify the item, gross price, applicant, utilisation, scientific need, available alternatives, market comparison, institutional commitments and required quotations, or use \string\DFGNotRequested{}.}

\DFGSubsection{5.2}{Module Temporary Position for Principal Investigator}

\DFGInstruction{Give the applicant-specific eligibility, duration, employment arrangement and scientific justification, or use \string\DFGNotRequested{}.}

\DFGSubsection{5.3}{Module Replacements}

\DFGInstruction{Give the applicant-specific eligibility, replacement arrangement, duration, calculation and justification, or use \string\DFGNotRequested{}.}

\DFGSubsection{5.4}{Module Temporary Substitutes for Clinicians}

\DFGInstruction{Give the clinical release arrangement, duration, calculation and justification, or use \string\DFGNotRequested{}.}

\DFGSubsection{5.5}{Module Mercator Fellows}

\DFGInstruction{Identify the fellow, contribution, visit plan, duration, calculation and added value, or use \string\DFGNotRequested{}.}

\DFGSubsection{5.6}{Module Project-Specific Workshops}

\DFGInstruction{Describe purpose, participants, programme, timing, calculation and project-specific necessity, or use \string\DFGNotRequested{}.}

\DFGSubsection{5.7}{Module Public Relations}

\DFGInstruction{Describe activities, audiences, timing, calculation and project relevance, or use \string\DFGNotRequested{}.}

\DFGSubsection{5.8}{Module Standard Allowance for Equity and Diversity}

\DFGInstruction{Describe eligible project-specific measures and the calculation, or use \string\DFGNotRequested{}. For Research Grants, check the current annual ceiling. For a network proposal, verify whether the request must be consolidated in the coordination project.}

% CONDITIONAL SERVICES MODULE
% Uncomment only when sequencing services exceed EUR 100,000 gross and are no
% more than EUR 1 million gross under form 54.020. In that case add the facility
% scientific director as an infrastructural co-applicant, that applicant's CV,
% the detailed cost estimate, and the section below.
%
% \DFGSubsection{5.9}{Services}
% \DFGInstruction{State the sequencing facility, infrastructural co-applicant,
% scope, sample-preparation and sequencing units, gross total, calculation,
% quality controls, responsibilities and required detailed cost estimate.}

% Attachments are uploaded separately in elan. Always include one current
% 53.200 CV per applicant. Add ethics statements, accepted manuscripts and
% editor confirmations, cooperation commitments, quotations, and other
% attachments only where the current instructions require them.

\end{document}
